\section{Recovering both arithmetic parameters from a finite signed zero list}
\label{sec:inverse}
The signed residues in the Lean proof retain surprisingly specific arithmetic
information. The earlier critical-line report \cite{priorzero} reconstructed the
modulus $q$ and auxiliary prime $d$ from their infinite sequence and asked for
a certified stopping rule with approximate data when both conductors are
unknown. We give such a rule in the parameter range of the lattice argument.
Only this section uses Dirichlet GRH as a premise.

\begin{theorem}\label{th:inverse}
Suppose $18480\mid q$, $d$ is prime and $d\equiv q/2-1\pmod q$.
In particular $d\equiv7\pmod8$, $(q,d)=1$, and $q\le2(d+1)$.
Set $h=\varphi(q)$, $\eta=(\cdot/d)$, and let $\chi^*$ denote the primitive
character inducing $\chi$. At $\rho=1/2+i\gamma$ the supplied ordinate weight is
\[
 m(\gamma)=h^{-1}\sum_{\chi\bmod q}
 \left(\operatorname{ord}_\rho L(s,\chi^*)-
       \operatorname{ord}_\rho L(s,\chi^*\eta)\right).
\]
Knowing the signed list through a sufficiently high finite height determines
$q,d$. The algorithm needs no prior numerical upper bound on either.
For approximate input it accepts enclosures for the finite sums it uses.
With complete lists and arbitrarily refinable certified errors it terminates.
For exact input a sufficient height is $O(d^2\log^3(2d))$ with an absolute
implied constant. This is a height estimate, independent of the cost of
computing or certifying the input lists.
\end{theorem}

\subsection*{Two terms in the explicit formula}
Write $\widehat g(\gamma)=\int g(t)e^{i\gamma t}\,dt$. For real compactly
supported test functions, the character-averaged Guinand--Weil formula takes
the form
\begin{align*}
 -\sum_\gamma m(\gamma)\widehat g(\gamma)
  ={}& (\log d)g(0)-\frac2h\int g(t)\cosh(t/2)\,dt\\
 &+\sum_{n=p^k}\frac{\log p}{\sqrt n}\,A(n)
                 [g(\log n)+g(-\log n)] .
\end{align*}
For $Q=p^{v_p(q)}$ the coefficient is
$A(n)=(1-\eta(n))\mathbf1_{n\equiv1\bmod(q/Q)}/\varphi(Q)$.
Here and below the sum over $n=p^k$ has $k\ge1$.
To check the identity, the primitive conductors in each pair are $f$ and $fd$,
so the signed minus sign yields $+\log d$ at the origin. Half of the characters
have each parity and twisting interchanges the two, canceling the gamma terms.
Only the untwisted family has the completed zeta poles. At $p\mid q$,
primitivization contributes only when the $p$-component is principal;
orthogonality over $q/Q$ gives the stated coefficient. Complex conjugates
occur in pairs. These are the terms of the classical explicit formula
\cite{chirre,booker}, specialized to the residue difference in
\path{New/Reduction/Zeros/PrimeError.lean}.

Consequently if $g$ kills the $\cosh(t/2)$ integral, has value one at $u$,
and its support isolates $u$, the measured quantity is $J(u)$, where
\[
 J(0)=\log d,\qquad J(\log p)=
 \frac{2\log p}{\sqrt p}\mathbf1_{q\mid p-1}
 \quad\text{for }p\nmid qd,\quad\eta(p)=-1 .
\]
The origin and prime identities alone would only give an asymptotic
reconstruction. To certify a finite stopping time we need to estimate the
unseen zero contribution before knowing the conductors.

\subsection*{A bound containing the unknown on both sides}
For a primitive character of conductor $f$, parity $a\in\{0,1\}$, let $B_f$
denote the sum of $(1+\gamma^2)^{-1}$ over its nontrivial zeros, including
multiplicity. Hadamard factorization and the functional equation give at $s=2$
\[
 \sum_\gamma\frac{3/2}{9/4+\gamma^2}=
 \Re\frac{L'}L(2)+\tfrac12\log(f/\pi)+\tfrac12\psi(1+a/2)
 +\tfrac32\mathbf1_{f=1}.
\]
This is the real logarithmic derivative of the completed entire function, with
$s(s-1)$ included for zeta \cite{lls}. GRH gives the indicated denominators.
Multiplying the summand by $3/2$ majorizes $(1+\gamma^2)^{-1}$.
For $f>1$, $|L'/L(2)|\le\sum_{n\ge2}(\log n)/n^2<1$ (isolate 2 and 3,
integrate from 3). For $f=1$ the logarithmic derivative is negative.
Since $\psi(1+a/2)\le\psi(2)<\log\pi$, in both cases
$B_f\le\tfrac34\log f+\tfrac94$.
Averaging absolute values before taking the signed difference yields
\[
 W:=\sum_\gamma\frac{|m(\gamma)|}{1+\gamma^2}
 \le \tfrac32\log q+\tfrac34\log d+\tfrac92
 \le \tfrac94\log d+7 ,
\]
using $q\le3d$. The dependence on $\log d$ here can be absorbed into the
origin measurement.

For any chosen test $g$, suppose $A$ bounds
$\sup_{|\gamma|>T}(1+\gamma^2)|\widehat g(\gamma)|$.
Let $[l,u]$ enclose the real part of the negative sum through height $T$.
The imaginary part of the exact full sum vanishes. For the origin probe put
$y=\log d$. Whenever $9A/4<1$, the arithmetic certificate is
\[
 \boxed{\quad \frac{l-7A}{1+9A/4}\ \le\ y\ \le \ \frac{u+7A}{1-9A/4}.\quad}
\]
Exponentiate with outward rounding. Stop the first stage when there is just
one admissible prime integer in the resulting interval. After that stage the
error for each new probe is at most $A(9\log d/4+7)$, independent of $q$.
A finite interval can likewise be certified by summing enclosures of
$m(\gamma)\widehat g(\gamma)$; omitted or uncertain entries require an enclosure
of their possible total contribution. Without such an input error guarantee a
list may omit a decisive term.

\subsection*{Two explicit probes}
A simple choice at location $v$ is the translate by $v$ of
\[
 g_a(t)=\frac{f_a''(t)-f_a(t)/4}{f_a''(0)-1/4},
 \qquad f_a(t)=\big(1-(t/a)^2\big)_+^8 .
\]
It equals one at its center and kills both $e^{t/2}$ and $e^{-t/2}$ by
integration by parts. Use $a=1/2$ at zero, or $a=1/[3(p+1)]$ at $\log p$;
the distance to the nearest other logarithm of a positive integer is at least
$\log(1+1/p)\ge1/(p+1)$. Four integrations by parts permit
$A=\|g_a^{(4)}\|_1(1+T^2)/T^4$. The accompanying script bounds this norm
term by term with rational arithmetic.

For the sharper height estimate, take $b=1/6$ at zero and
$b=1/[10(p+1)]$ at $\log p$. Set $k=\lfloor bT/2\rfloor\ge4$.
Convolve the uniform probability density on $[-b/k,b/k]$ with itself $k$
times and divide by its value at zero to obtain $w$. It is symmetric and
unimodal (averaging any symmetric unimodal density on a centered interval
preserves this property). Before normalization its maximum is at least $1/(2b)$
by total mass and support. Define
\[
 g(t)=w(t-v)-\frac{w(t-v-2b)+w(t-v+2b)}{2\cosh b}.
\]
The support lies within $3b$ of the center, $g(v)=1$, and translation of
$e^{\pm t/2}$ by $\pm2b$ cancels its integral. The Fourier transform of the
unnormalized $w$ is $\operatorname{sinc}(b\gamma/k)^k$, where
$\operatorname{sinc}(x)=\sin(x)/x$. For $|\gamma|>T$ the absolute
value is at most $2^{-k}(T/|\gamma|)^k$. The translation multiplier is at
most two, so we can use
\[
 A=4b(1+T^2)\,2^{-\lfloor bT/2\rfloor}.
\]
The probes have sufficient differentiability for the explicit formula:
convolution with smooth approximate identities gives the same identity in
the limit, with spectral domination by a multiple of $(1+\gamma^2)^{-1}$.
Increasing $T$ and refining enclosures therefore certifies $d$ at the
origin. One can, for example, choose $T$ dyadically and refine the finite
sum enclosure to width $1/T^2$ or smaller; adaptive refinement to arbitrary
accuracy at each trial also works. For exact data $T=O(\log(2d))$ suffices at
the origin since adjacent integer logarithms near $d$ are separated by
$\gg1/d$. This last scale estimate assumes exact finite sums.

\subsection*{Separating the finitely many moduli}
Form $M=2(d+1)$. The remaining candidates are the divisors $q$ of $M$
with $18480\mid q$ and $M/q$ odd. For two distinct candidates $q_1<q_2$,
there is a unit $u\bmod M$ congruent to one modulo $q_1$ and different
from one modulo $q_2$. Indeed choose a prime $\ell$ with
$v_\ell(q_1)<v_\ell(q_2)$. Prescribe the residue $1+\ell^{v_\ell(q_1)}$
modulo the full $\ell$-power in $M$, and one at all other prime powers.
If $v_\ell(q_1)=0$, this residue is two and $\ell$ is odd. Now require
$p\equiv u\bmod M$ and $\eta(p)=-1$. These conditions specify one coset of
index $2\varphi(M)$ modulo $Md$, since $(M,d)=1$. Dirichlet's theorem
gives a prime in this coset, and a search with integer tests finds one.
The two possible values of $J(\log p)$ are separated by $2\log p/\sqrt p$.
A certified interval excluding one value eliminates a candidate. Continuing
until only one remains proves finite recovery.

Lamzouri--Li--Soundararajan's Theorem 1.4 bounds the least prime in a coset
of index $I>1$ modulo $Q\ge20000$ by $O(I^2\log^2 Q)$ under Dirichlet GRH
\cite{lls}. Their explicit alternatives are $p\le10^9$ or
$p\le((I-1)\log Q+3(I+1)+\tfrac52(\log\log Q)^2)^2$. Here $I\le 2M$ and $Q=Md$, so each
witness can be taken $p\ll d^2\log^2(2d)$. For a sufficiently large absolute
constant, $T$ of order $p\log(2pd)$ makes the spline tail times $W$ less than
half the two-value separation. All witnesses use the same spectral list, so
the maximum of these heights proves the claimed bound.
\hfill$\square$

\paragraph{An arithmetic example.}
Take $d=109909799$. Sixteen candidate moduli survive the divisor tests. Even
knowing the totient would leave the pair $1127280,1201200$ at totient
$230400$. The prime $34945681$ is a quadratic nonresidue modulo $d$, and
$\gcd(p-1,2(d+1))=1127280$; its jump distinguishes the two. At $T=768$ the
spline at the origin has a tail below $1.045\cdot10^{-12}$ for this actual
$d$. At $T=6\cdot10^{10}$ the witness probe has a tail below
$5.220\cdot10^{-11}$, against a jump greater than $0.005875$.
The script checks primes, characters, candidate enumeration and tail
arithmetic. These are calculations on known parameters; an observed
spectral list of these heights was not computed.
